\documentclass[10pt,aspectratio=169]{beamer}

% ============================================================
% Encodage et langue
% ============================================================
\usepackage[utf8]{inputenc}
\usepackage[T1]{fontenc}
\usepackage[french]{babel}
\usepackage{lmodern}
\usepackage{microtype}

% ============================================================
% Packages graphiques
% ============================================================
\usepackage{tikz}
\usetikzlibrary{arrows.meta,positioning,shapes.geometric}
\usepackage{booktabs}
\usepackage{array}
\usepackage{xcolor}
\usepackage{listings}

% ============================================================
% Style
% ============================================================
\definecolor{MainBlue}{HTML}{0B5C8E}
\definecolor{DarkBlue}{HTML}{16324F}
\definecolor{LightBox}{HTML}{E9EEF3}
\definecolor{VeryLight}{HTML}{F6F8FA}
\definecolor{TextGray}{HTML}{263445}

\setbeamertemplate{navigation symbols}{}
\setbeamertemplate{footline}[frame number]
\setbeamercolor{title}{fg=MainBlue}
\setbeamercolor{frametitle}{fg=MainBlue}
\setbeamerfont{frametitle}{series=\bfseries,size=\Large}
\setbeamercolor{normal text}{fg=TextGray}

\setlength{\parskip}{0.25em}

\newcommand{\smallbullet}{\textcolor{MainBlue}{$\blacktriangleright$}\hspace{0.4em}}

\newenvironment{infobox}[1]
{\begin{beamercolorbox}[rounded=true,shadow=false,wd=\linewidth,sep=0.8em]{block body}
\textbf{\color{DarkBlue}#1}\par\vspace{0.3em}}
{\end{beamercolorbox}}

\setbeamercolor{block body}{bg=LightBox,fg=TextGray}

\lstset{
  basicstyle=\ttfamily\small,
  columns=fullflexible,
  keepspaces=true,
  frame=none,
  backgroundcolor=\color{VeryLight},
  xleftmargin=0.5em,
  xrightmargin=0.5em
}

% ============================================================
% Métadonnées
% ============================================================
\title{Scripts d'extractions et de traitement de données de vie réelle d'ARIA à des fins de recherche clinique}
\subtitle{ARIA ODM Builder\\Notice visuelle du dépôt GitHub}
\author{}
\date{}

\begin{document}

% ============================================================
\begin{frame}
    \titlepage
    \vspace{-1.2em}
    \centering
    \begin{tikzpicture}[node distance=0.25cm, every node/.style={font=\small}]
        \node[draw=MainBlue, rounded corners, minimum width=2.3cm, minimum height=0.75cm] (sql) {SQL};
        \node[draw=MainBlue, rounded corners, minimum width=2.3cm, minimum height=0.75cm, right=of sql] (script) {Script principal};
        \node[draw=MainBlue, rounded corners, minimum width=2.3cm, minimum height=0.75cm, right=of script] (files) {CSV/XLSX};
        \node[draw=MainBlue, rounded corners, minimum width=2.3cm, minimum height=0.75cm, right=of files] (streamlit) {Streamlit};
        \node[draw=MainBlue, rounded corners, minimum width=2.3cm, minimum height=0.75cm, right=of streamlit] (export) {Export};
        \draw[-{Latex}, thick, MainBlue] (sql) -- (script);
        \draw[-{Latex}, thick, MainBlue] (script) -- (files);
        \draw[-{Latex}, thick, MainBlue] (files) -- (streamlit);
        \draw[-{Latex}, thick, MainBlue] (streamlit) -- (export);
    \end{tikzpicture}
\end{frame}

% ============================================================
\begin{frame}{Comment lire ce document}
\begin{columns}[T,onlytextwidth]
\begin{column}{0.48\linewidth}
\begin{infobox}{But du PDF}
Ce document est conçu pour être lu directement depuis le dossier \texttt{docs/} du dépôt GitHub.

\vspace{0.5em}
\smallbullet comprendre l'ordre réel du pipeline ;\\
\smallbullet identifier les fichiers importants ;\\
\smallbullet savoir quel script sert à quoi ;\\
\smallbullet comprendre le résultat attendu.
\end{infobox}
\end{column}

\begin{column}{0.48\linewidth}
\begin{infobox}{Public visé}
Utilisateur non spécialiste du code : médecin, physicien médical, interne, chercheur ou collaborateur projet.

\vspace{0.7em}
Ce PDF fonctionne comme un README visuel.
\end{infobox}
\end{column}
\end{columns}
\end{frame}

% ============================================================
\begin{frame}{Objectif du dépôt}
\centering
\Large
Transformer des extractions ARIA en fichiers patients structurés, contrôlés et exploitables pour la recherche clinique.

\vspace{1.1cm}

\begin{columns}[T,onlytextwidth]
\begin{column}{0.31\linewidth}
\begin{infobox}{Extraire}
Récupérer les informations utiles avec les requêtes SQL.
\end{infobox}
\end{column}

\begin{column}{0.31\linewidth}
\begin{infobox}{Préparer}
Générer les fichiers d'entrée CSV et Excel avec le script principal.
\end{infobox}
\end{column}

\begin{column}{0.31\linewidth}
\begin{infobox}{Exporter}
Produire un résultat structuré avec l'application Streamlit.
\end{infobox}
\end{column}
\end{columns}
\end{frame}

% ============================================================
\begin{frame}{Chaîne complète du projet}
\centering
\begin{tikzpicture}[
    node distance=0.65cm,
    step/.style={draw=MainBlue, fill=LightBox, rounded corners, minimum width=2.75cm, minimum height=0.9cm, align=center, font=\small\bfseries}
]
\node[step] (sql) {1. \texttt{sql/}};
\node[step, right=of sql] (load) {2. \texttt{load\_final\_all.py}};
\node[step, right=of load] (files) {3. CSV/XLSX};
\node[step, below=1.0cm of files] (conf) {4. \texttt{conf/}};
\node[step, left=of conf] (app) {5. \texttt{app.py}};
\node[step, left=of app] (res) {6. Résultats};

\draw[-{Latex}, thick, MainBlue] (sql) -- (load);
\draw[-{Latex}, thick, MainBlue] (load) -- (files);
\draw[-{Latex}, thick, MainBlue] (files) -- (conf);
\draw[-{Latex}, thick, MainBlue] (conf) -- (app);
\draw[-{Latex}, thick, MainBlue] (app) -- (res);
\end{tikzpicture}

\vspace{0.85cm}
\begin{flushleft}
\smallbullet \texttt{sql/} contient les requêtes d'extraction.\\
\smallbullet \texttt{scripts/load\_final\_all.py} lance les requêtes et génère les fichiers d'entrée.\\
\smallbullet \texttt{app.py} charge, contrôle et exporte les données via Streamlit.
\end{flushleft}
\end{frame}

% ============================================================
\begin{frame}[fragile]{Organisation visible sur GitHub}
\begin{columns}[T,onlytextwidth]
\begin{column}{0.46\linewidth}
\begin{lstlisting}
ARIA_ODM_Builder/
|-- conf/
|-- docs/
|-- pictures/
|-- samples/
|-- scripts/
|-- sql/
|-- utils/
|-- README.md
|-- app.py
|-- requirements.txt
`-- requirements_sql.txt
\end{lstlisting}
\end{column}

\begin{column}{0.50\linewidth}
\begin{infobox}{Lecture rapide}
Le dépôt suit la chaîne : extraction SQL, préparation par scripts, fichiers d'entrée, configuration, application Streamlit et export.
\end{infobox}

\vspace{0.5cm}

\begin{infobox}{Point d'entrée utilisateur}
Pour l'utilisation guidée, l'utilisateur lance \texttt{app.py}. Pour la préparation amont, le script principal est dans \texttt{scripts/}.
\end{infobox}
\end{column}
\end{columns}
\end{frame}

% ============================================================
\begin{frame}[fragile]{Étape 1 : requêtes SQL dans \texttt{sql/}}
\begin{columns}[T,onlytextwidth]
\begin{column}{0.46\linewidth}
\begin{lstlisting}
sql/
|-- README.md
|-- requete_traitement.sql
|-- requete_formulaire.sql
`-- requete_complementaire.sql
\end{lstlisting}
\end{column}

\begin{column}{0.50\linewidth}
\begin{infobox}{Rôle}
Produire les premières extractions nécessaires au pipeline.

\vspace{0.4em}
\smallbullet extraction des données de traitement ;\\
\smallbullet extraction des formulaires cliniques ;\\
\smallbullet extraction d'une source complémentaire si nécessaire.
\end{infobox}
\end{column}
\end{columns}

\vfill
\centering
\textbf{Guide détaillé :} \texttt{docs/guide\_execution\_sql.md}
\end{frame}

% ============================================================
\begin{frame}[fragile]{Connexion à la base ARIA : où renseigner les paramètres}
\begin{columns}[T,onlytextwidth]
\begin{column}{0.46\linewidth}
\begin{infobox}{Où modifier ?}
Les informations de connexion sont à renseigner dans le script principal de préparation :

\vspace{0.4em}
\texttt{scripts/load\_final\_all.py}

\vspace{0.4em}
fonction : \texttt{connecter\_bdd()}\\
chercher : \texttt{pyodbc.connect(...)}

\vspace{0.4em}
\smallbullet nom du serveur ;\\
\smallbullet nom de la base ;\\
\smallbullet mode d'authentification ;\\
\smallbullet port si nécessaire.
\end{infobox}
\end{column}

\begin{column}{0.50\linewidth}
\begin{infobox}{Exemple à adapter}
\begin{lstlisting}
pyodbc.connect(
"Driver={ODBC Driver 17 for SQL Server};"
"Server=REMPLIR_ICI_NOM_SERVEUR,PORT;"
"Database=REMPLIR_ICI_NOM_BASE;"
"Trusted_Connection=yes;"
"TrustServerCertificate=yes;"
)
\end{lstlisting}
\end{infobox}
\end{column}
\end{columns}

\vfill
\centering
\textbf{À retenir :} l'utilisateur complète ce bloc avec les paramètres fournis par son environnement de travail.
\end{frame}

% ============================================================
\begin{frame}[fragile]{Étape 2 : préparation dans \texttt{scripts/}}
\begin{columns}[T,onlytextwidth]
\begin{column}{0.44\linewidth}
\begin{lstlisting}
scripts/
|-- README.md
`-- load_final_all.py
\end{lstlisting}
\end{column}

\begin{column}{0.52\linewidth}
\begin{infobox}{Script principal}
\texttt{load\_final\_all.py} centralise l'extraction amont.

\vspace{0.5em}
Il exécute les requêtes SQL, récupère les résultats et génère les fichiers d'entrée au format CSV et Excel.
\end{infobox}
\end{column}
\end{columns}
\end{frame}

% ============================================================
\begin{frame}{Sortie de la préparation amont}
\begin{columns}[T,onlytextwidth]
\begin{column}{0.50\linewidth}
\centering
\begin{tikzpicture}[
    node distance=0.42cm,
    box/.style={draw=MainBlue, fill=LightBox, rounded corners, text width=3.8cm, align=center, minimum height=0.72cm}
]
\node[box] (sql) {Requêtes SQL};
\node[box, below=of sql] (load) {\texttt{load\_final\_all.py}};
\node[box, below=of load] (csv) {Fichiers CSV};
\node[box, below=of csv] (xlsx) {Fichiers Excel};
\node[box, below=of xlsx] (app) {Chargement dans Streamlit};

\draw[-{Latex}, thick, MainBlue] (sql) -- (load);
\draw[-{Latex}, thick, MainBlue] (load) -- (csv);
\draw[-{Latex}, thick, MainBlue] (csv) -- (xlsx);
\draw[-{Latex}, thick, MainBlue] (xlsx) -- (app);
\end{tikzpicture}
\end{column}

\begin{column}{0.46\linewidth}
\begin{infobox}{Fichiers obtenus}
Le script produit les fichiers nécessaires au workflow Streamlit dans deux formats : \texttt{.csv} et \texttt{.xlsx}.
\end{infobox}

\vspace{0.6cm}

\begin{infobox}{Intérêt}
Cette étape évite de mélanger extraction, génération des fichiers d'entrée et utilisation interactive.
\end{infobox}
\end{column}
\end{columns}
\end{frame}

% ============================================================
\begin{frame}[fragile]{Étape 3 : fichiers d'entrée pour Streamlit}
\begin{columns}[T,onlytextwidth]
\begin{column}{0.48\linewidth}
\begin{infobox}{Cas réel : fichiers générés}
\begin{lstlisting}
traitement_patient.csv / .xlsx
formulaire_patient.csv / .xlsx
ethos_patient.csv / .xlsx
\end{lstlisting}
\end{infobox}
\end{column}

\begin{column}{0.48\linewidth}
\begin{infobox}{Cas démo : fichiers fictifs}
\begin{lstlisting}
samples/
|-- csv/
`-- xlsx/
\end{lstlisting}
Les fichiers de démonstration permettent de tester l'application sans extraction SQL.
\end{infobox}
\end{column}
\end{columns}

\vfill
\centering
\textbf{Ces fichiers sont ceux que l'utilisateur charge ensuite dans l'interface Streamlit.}
\end{frame}

% ============================================================
\begin{frame}{Étape 4 : configuration dans \texttt{conf/}}
\begin{columns}[T,onlytextwidth]
\begin{column}{0.48\linewidth}
\begin{infobox}{\texttt{mapping.csv}}
Indique quelle colonne source correspond à quelle variable utilisée dans le pipeline.

\vspace{0.4em}
Exemples : identifiant patient, date, dose, diagnostic, fractionnement.
\end{infobox}
\end{column}

\begin{column}{0.48\linewidth}
\begin{infobox}{\texttt{profiles/*.json}}
Décrit un scénario d'extraction : pathologie, critères, variables attendues et règles de sélection.

\vspace{0.4em}
Exemples : codes CIM10, variables cliniques, règles de cohorte.
\end{infobox}
\end{column}
\end{columns}

\vfill
\centering
\textbf{Le mapping dit où lire. Le profil dit quoi sélectionner.}
\end{frame}

% ============================================================
\begin{frame}[fragile]{Installer le projet}
\begin{columns}[T,onlytextwidth]
\begin{column}{0.48\linewidth}
\begin{infobox}{Application Streamlit}
\begin{lstlisting}
python -m venv .venv
.venv\Scripts\activate
pip install -r requirements.txt
\end{lstlisting}
\texttt{requirements.txt} contient les dépendances de l'application.
\end{infobox}
\end{column}

\begin{column}{0.48\linewidth}
\begin{infobox}{Extraction SQL}
\begin{lstlisting}
python -m venv .venv_sql
.venv_sql\Scripts\activate
pip install -r requirements_sql.txt
\end{lstlisting}
\texttt{requirements\_sql.txt} contient les dépendances minimales pour l'extraction SQL.
\end{infobox}
\end{column}
\end{columns}
\end{frame}

% ============================================================
\begin{frame}[fragile]{Étape 5 : lancer l'application Streamlit}
\begin{columns}[T,onlytextwidth]
\begin{column}{0.46\linewidth}
\begin{infobox}{Commande}
\begin{lstlisting}
python -m streamlit cache clear
python -m streamlit run app.py
\end{lstlisting}

\smallbullet l'application s'ouvre dans le navigateur ;\\
\smallbullet l'utilisateur charge les fichiers préparés ;\\
\smallbullet le pipeline peut être recalculé depuis l'interface.
\end{infobox}
\end{column}

\begin{column}{0.50\linewidth}
\centering
\begin{tikzpicture}[
    node distance=0.7cm,
    box/.style={draw=MainBlue, fill=LightBox, rounded corners, text width=4.2cm, align=center, minimum height=0.85cm}
]
\node[box] (terminal) {Terminal\\\texttt{python -m streamlit run app.py}};
\node[box, below=of terminal] (browser) {Interface dans le navigateur};
\draw[-{Latex}, thick, MainBlue] (terminal) -- (browser);
\end{tikzpicture}
\end{column}
\end{columns}
\end{frame}

% ============================================================
\begin{frame}{Ce que fait \texttt{app.py}}
\centering
\begin{tikzpicture}[
    node distance=0.35cm,
    box/.style={draw=MainBlue, fill=LightBox, rounded corners, minimum width=1.9cm, minimum height=0.75cm, align=center, font=\small}
]
\node[box] (a) {Charger};
\node[box, right=of a] (b) {Mapper};
\node[box, right=of b] (c) {Filtrer};
\node[box, right=of c] (d) {Joindre};
\node[box, right=of d] (e) {Contrôler};
\node[box, right=of e] (f) {Exporter};

\foreach \x/\y in {a/b,b/c,c/d,d/e,e/f}
    \draw[-{Latex}, thick, MainBlue] (\x) -- (\y);
\end{tikzpicture}

\vspace{1cm}
\begin{flushleft}
\smallbullet interface principale du dépôt ;\\
\smallbullet guide l'utilisateur étape par étape ;\\
\smallbullet s'appuie sur les modules du dossier \texttt{utils/}.
\end{flushleft}

\vfill
\begin{infobox}{À retenir}
L'utilisateur standard lance \texttt{app.py} après la préparation des fichiers d'entrée.
\end{infobox}
\end{frame}

% ============================================================
\begin{frame}{Modules Python utilisés par l'application : \texttt{utils/}}
\small
\begin{tabular}{>{\ttfamily}p{0.22\linewidth} p{0.70\linewidth}}
\toprule
\textbf{Module} & \textbf{Rôle}\\
\midrule
load.py & charger les fichiers d'entrée\\
clean.py & nettoyer et standardiser les données utiles\\
mapping.py & appliquer les correspondances de variables\\
cohort.py & construire la cohorte selon les critères choisis\\
temporal.py & calculer les variables liées au temps\\
quality.py & produire les contrôles qualité\\
export.py & générer les fichiers de sortie\\
profile.py & lire et valider les profils JSON\\
display.py & afficher les résultats dans Streamlit\\
text.py & gérer les textes, libellés et messages\\
\bottomrule
\end{tabular}
\end{frame}

% ============================================================
\begin{frame}{Documentation disponible dans \texttt{docs/}}
\small
\begin{tabular}{>{\ttfamily}p{0.34\linewidth} p{0.58\linewidth}}
\toprule
\textbf{Fichier} & \textbf{Rôle}\\
\midrule
overview\_github.pdf & notice visuelle du dépôt GitHub\\
overview\_github.tex & source LaTeX modifiable de cette notice\\
documentation.tex & documentation utilisateur détaillée\\
guide\_execution\_sql.md & guide pour lancer les extractions SQL\\
scenarios.md & scénarios d'utilisation de l'application\\
scenarios.tex & source LaTeX des scénarios\\
README.md & sommaire du dossier \texttt{docs/}\\
\bottomrule
\end{tabular}
\end{frame}

% ============================================================
\begin{frame}{Résultats attendus après exécution}
\begin{columns}[T,onlytextwidth]
\begin{column}{0.46\linewidth}
\begin{infobox}{Fichiers produits}
\smallbullet export patient structuré ;\\
\smallbullet rapport qualité ;\\
\smallbullet trace d'exécution.
\end{infobox}
\end{column}

\begin{column}{0.50\linewidth}
\begin{infobox}{Utilisation}
Ces sorties servent de base au contrôle, à la relecture puis aux analyses de recherche clinique.

\vspace{0.4em}
\smallbullet vérifier la cohorte ;\\
\smallbullet relire les alertes qualité ;\\
\smallbullet exploiter l'export final.
\end{infobox}
\end{column}
\end{columns}

\vfill
\centering
\textbf{L'utilisateur sait alors quels fichiers ont été utilisés et quel résultat a été généré.}
\end{frame}

% ============================================================
\begin{frame}{Workflow résumé}
\centering
\begin{tikzpicture}[
    node distance=0.35cm,
    box/.style={draw=MainBlue, fill=LightBox, rounded corners, minimum width=1.85cm, minimum height=0.75cm, align=center, font=\small}
]
\node[box] (s1) {1. SQL};
\node[box, right=of s1] (s2) {2. Script};
\node[box, right=of s2] (s3) {3. CSV/XLSX};
\node[box, right=of s3] (s4) {4. Profil};
\node[box, right=of s4] (s5) {5. App};
\node[box, right=of s5] (s6) {6. Export};

\foreach \x/\y in {s1/s2,s2/s3,s3/s4,s4/s5,s5/s6}
    \draw[-{Latex}, thick, MainBlue] (\x) -- (\y);
\end{tikzpicture}

\vspace{0.8cm}

\begin{flushleft}
\begin{enumerate}
    \item Exécuter les requêtes du dossier \texttt{sql/}.
    \item Générer les fichiers avec \texttt{scripts/load\_final\_all.py}.
    \item Charger les fichiers CSV/XLSX dans Streamlit ou utiliser \texttt{samples/}.
    \item Choisir le mapping et le profil dans \texttt{conf/}.
    \item Lancer les contrôles qualité.
    \item Générer l'export final.
\end{enumerate}
\end{flushleft}
\end{frame}

\end{document}
